\documentclass{article}
\usepackage[utf8]{inputenc}
\usepackage{amsmath,amssymb}
\usepackage[most]{tcolorbox}
\usepackage{geometry}
\geometry{margin=2.5cm}

\newcommand{\step}[1]{\par\noindent\hangindent=3.5em\hangafter=1%
  \makebox[3.5em][l]{\textbf{#1.}}}
\newcommand{\steptag}[1]{\hfill{\small\textsf{[#1]}}}

\newtcolorbox{proofbox}[1]{
  colback=gray!5, colframe=gray!60,
  fonttitle=\bfseries, title={#1},
  breakable, enhanced,
  left=4pt, right=4pt, top=4pt, bottom=4pt,
  before skip=10pt, after skip=10pt
}

\begin{document}

\begin{proofbox}{Explicit group-input polar-domain membership: there exist...}
\step{1} Explicit group-input polar-domain membership: there exist universal C\_grp, C\_pol >= 1, kappa\_pol in (0, 1/2] such that for every finite-dimensional exact-unit epsilon\_r-C*-algebra and delta > 0 satisfying C\_pol*(epsilon\_r + delta) <= kappa\_pol and C\_grp*epsilon\_r < delta - C\_pol*(epsilon\_r*delta + delta\textasciicircum{}2), writing (u\_delta, h\_delta) for the unique inverse of Pi\_delta: calU x B\textasciicircum{}\{calH\}\_delta(J) -> S\_delta := Pi\_delta(calU x B\textasciicircum{}\{calH\}\_delta(J)), Pi\_delta(U, H) = U bold-dot H, the first inverse component u\_delta:S\_delta -> calU is defined at U bold-dot V and U\textasciicircum{}dagger for every U, V in calU; moreover, U bold-dot V and U\textasciicircum{}dagger each have a right inverse. \steptag{validated} \\[6pt]
\step{1.1} Invoke the registered external lem-stage1-polar-retraction and choose its universal constants C\_pol >= 1 and kappa\_pol in (0,1/2]; choose the further universal witness C\_grp:=4. These constants are fixed for the rest of the proof. \steptag{validated} \\[4pt]
\step{1.2} Fix an arbitrary finite-dimensional exact-unit epsilon\_r-C*-algebra and delta>0 satisfying the two root guards. Apply lem-stage1-polar-retraction under the first guard to obtain the displayed C\textasciicircum{}1 diffeomorphism Pi\_delta:calU x B\_delta\textasciicircum{}\{calH\}(J)->S\_delta and, before choosing any group inputs, bind its unique typed inverse pair (u\_delta,h\_delta):S\_delta->calU x B\_delta\textasciicircum{}\{calH\}(J). Put e:=epsilon\_r and t:=delta-C\_pol*(e*delta+delta\textasciicircum{}2). The external gives the inner inclusion calU\_t subseteq S\_delta. \steptag{validated} \\[4pt]
\step{1.3} Now fix arbitrary U,V in calU. Guard arithmetic gives s:=e+delta<=1/2, t>4e>=0, delta<=1/2, and e<1/6: indeed C\_pol*s<=kappa\_pol<=1/2 with C\_pol>=1 gives s<=1/2, while the second guard gives 4e<t; also t=delta-C\_pol*delta*s<=delta*(1-s)=(s-e)*(1-s), whose maximum for e<=s<=1/2 is (1/2-e)/2, so e<t<=(1/2-e)/2 implies e<1/6. \steptag{validated} \\[4pt]
\step{1.4} Common left-multiplier estimate, using the preceding guard node: for every W in calU, def-approximate-unitary-space gives W\textasciicircum{}dagger bold-dot W=J and a right inverse, while def-epsilon-cstar-algebra gives (1-e)||W||\textasciicircum{}2<=1 and ||L\_\{W\textasciicircum{}dagger\}L\_W-I||<=e||W||\textasciicircum{}2<=e/(1-e)<1/5. Hence L\_W is injective and, since calX is finite-dimensional, bijective. Moreover ||W||<=6/5, ||L\_\{W\textasciicircum{}dagger\}||<=(1+e)||W||<=7/5, and ||L\_\{W\textasciicircum{}dagger\}L\_W Z||>=(4/5)||Z||, so ||L\_W Z||>=(4/7)||Z|| and ||L\_W\textasciicircum{}\{-1\}||<=7/4. \steptag{validated} \\[4pt]
\step{1.5} Product defect: put X:=U bold-dot V. The star rule gives X\textasciicircum{}dagger=V\textasciicircum{}dagger bold-dot U\textasciicircum{}dagger. Two applications of the approximate-associativity axiom, inserting ((V\textasciicircum{}dagger bold-dot U\textasciicircum{}dagger) bold-dot U) bold-dot V and (V\textasciicircum{}dagger bold-dot (U\textasciicircum{}dagger bold-dot U)) bold-dot V, and then U\textasciicircum{}dagger bold-dot U=V\textasciicircum{}dagger bold-dot V=J, give ||X\textasciicircum{}dagger bold-dot X-J||<=2e(1+e)||U||\textasciicircum{}2||V||\textasciicircum{}2<=2e(1+e)/(1-e)\textasciicircum{}2<=4e<t<2t. This uses only the registered definitions and the guard estimates above. \steptag{validated} \\[4pt]
\step{1.6} Product right inverse: approximate associativity gives ||L\_X-L\_U L\_V||<=e||U||||V||<=e/(1-e)<1/5. The operator A:=L\_U L\_V is invertible by the common multiplier node and ||A\textasciicircum{}\{-1\}||<=49/16, so ||A\textasciicircum{}\{-1\}(L\_X-A)||<49/80<1. The Neumann perturbation lemma makes L\_X invertible; therefore R\_X:=L\_X\textasciicircum{}\{-1\}J satisfies X bold-dot R\_X=J and is a right inverse of X. \steptag{validated} \\[6pt]
\step{1.6.1} Bind e:=epsilon\_r and X:=U bold-dot V locally, with U,V in calU fixed by the parent context. The root guards imply e<1/6: setting s:=e+delta and t:=delta-C\_pol*delta*s, one has s<=1/2 and 4e<t<=delta*(1-s)=(s-e)*(1-s)<=(1/2-e)/2, hence e<1/6. For any W in calU, W\textasciicircum{}dagger bold-dot W=J by def-approximate-unitary-space, while the C*-lower bound and exact unit give (1-e)||W||\textasciicircum{}2<=||W\textasciicircum{}dagger bold-dot W||=1, so ||W||\textasciicircum{}2<=1/(1-e) and ||W||<6/5. Approximate associativity gives ||L\_\{W\textasciicircum{}dagger\}L\_W-I||<=e||W||\textasciicircum{}2<=e/(1-e)<1/5. Thus L\_W is injective (if L\_W Z=0, then ||Z||<=(e/(1-e))||Z||), hence bijective in finite dimension. Also ||L\_\{W\textasciicircum{}dagger\}||<=(1+e)||W||<7/5 and ||L\_\{W\textasciicircum{}dagger\}L\_W Z||>(4/5)||Z||, whence ||L\_W Z||>(4/7)||Z|| and ||L\_W\textasciicircum{}\{-1\}||<=7/4. Applying this self-contained derivation to W=U,V, A:=L\_U L\_V is invertible with A\textasciicircum{}\{-1\}=L\_V\textasciicircum{}\{-1\}L\_U\textasciicircum{}\{-1\} and ||A\textasciicircum{}\{-1\}||<=49/16. Approximate associativity and ||U||||V||<=1/(1-e) give ||L\_X-A||<=e||U||||V||<=e/(1-e)<1/5. Therefore ||A\textasciicircum{}\{-1\}(L\_X-A)||<49/80<1. By the Neumann lemma, A\textasciicircum{}\{-1\}L\_X=I+A\textasciicircum{}\{-1\}(L\_X-A), hence L\_X, is invertible. Consequently R\_X:=L\_X\textasciicircum{}\{-1\}J satisfies X bold-dot R\_X=L\_X(R\_X)=J, so X has the required right inverse. \steptag{validated} \\[4pt]
\step{1.7} Adjoint defect and right inverse: def-approximate-unitary-space gives U\textasciicircum{}dagger bold-dot U=J, so U is already a right inverse of U\textasciicircum{}dagger. Let R be the right inverse of U supplied by that definition. Bijectivity of L\_U from the common multiplier node gives R=L\_U\textasciicircum{}\{-1\}J and ||R||<=7/4. Exact unitality and one associator estimate yield ||U\textasciicircum{}dagger-R||=||U\textasciicircum{}dagger bold-dot(U bold-dot R)-(U\textasciicircum{}dagger bold-dot U) bold-dot R||<=e||U||\textasciicircum{}2||R||. Hence ||U bold-dot U\textasciicircum{}dagger-J||<= (1+e)e||U||\textasciicircum{}3||R||<=4e<t<2t, using e<1/6, ||U||<=6/5, and (7/6)*(6/5)\textasciicircum{}3*(7/4)=441/125<4. \steptag{validated} \\[4pt]
\step{1.8} The product-defect and product-right-inverse nodes show X=U bold-dot V lies in calU\_t by def-approximate-unitary-space. The adjoint node shows U\textasciicircum{}dagger lies in calU\_t, because its defining defect is ||(U\textasciicircum{}dagger)\textasciicircum{}dagger bold-dot U\textasciicircum{}dagger-J||=||U bold-dot U\textasciicircum{}dagger-J||<2t and U is its right inverse. Thus both required right inverses are also explicitly recorded. \steptag{validated} \\[4pt]
\step{1.9} By the typed polar data bound before the group inputs, lem-stage1-polar-retraction supplies calU\_t subseteq S\_delta. The preceding membership node therefore puts U bold-dot V and U\textasciicircum{}dagger in the domain S\_delta of the already bound map u\_delta:S\_delta->calU, so u\_delta is defined at both points. Since the algebra, delta, U, and V were arbitrary, C\_grp=4 together with the external C\_pol,kappa\_pol proves every clause of the root contract. \steptag{validated} \\[6pt]
\step{1.9.1} By validated node 1.2, for t:=delta-C\_pol*(epsilon\_r*delta+delta\textasciicircum{}2), the already bound inverse has first component u\_delta:S\_delta->calU and the polar external gives calU\_t subseteq S\_delta. By validated node 1.8, U bold-dot V and U\textasciicircum{}dagger both lie in calU\_t, and it also establishes that each has a right inverse. Hence the inclusion places both points in S\_delta, so the fixed map u\_delta is defined at both. Together with the universal choices C\_grp=4 and C\_pol,kappa\_pol from validated node 1.1, and the arbitrary algebra, delta, U,V fixed in validated nodes 1.2-1.3, this proves every clause of the root contract. \steptag{archived} \\[4pt]
\step{1.9.2} By node 1.2, with t:=delta-C\_pol*(epsilon\_r*delta+delta\textasciicircum{}2), the inverse pair was bound before U,V and its first component is u\_delta:S\_delta->calU, while calU\_t subseteq S\_delta. Node 1.8 proves U bold-dot V in calU\_t and U\textasciicircum{}dagger in calU\_t and proves that each point has a right inverse. Therefore U bold-dot V and U\textasciicircum{}dagger belong to S\_delta, so the already bound u\_delta is defined at both. Node 1.1 supplies the universal choices C\_grp=4 and C\_pol,kappa\_pol, and nodes 1.2-1.3 supply the arbitrary guarded algebra, delta, U,V; universal generalization now yields every clause of the root contract. \steptag{validated} \\[4pt]
\end{proofbox}

\end{document}
